\documentclass[10pt,a4paper]{amsart}
\usepackage[margin=19mm]{geometry}
\usepackage{amsmath,amssymb,mathtools}
\usepackage[scr=rsfs,cal=euler]{mathalfa}
\usepackage{xfrac}
\usepackage[unicode,hidelinks,bookmarks=false]{hyperref}
\newcommand{\ie}{i.e.\ }
\newcommand{\cf}{cf.\ }
\newcommand{\resp}{respectively\ }
\newcommand{\Subsection}{\subsection}
\providecommand{\nobreakdash}{\nobreak\hbox{-}}
\setlength{\emergencystretch}{2em}
\multlinegap0pt
\usepackage{amssymb}
\usepackage{mathtools}
\usepackage{xcolor}
\newtheorem{prop}{Proposition}
\title{Fenchel-Willmore-Chen Inequality under Lower Bounds on weighted intermediate Ricci curvature}
\author{Jihye Lee, Guofang Wei}
\date{Source: arXiv:2608.01171v1 (CC BY 4.0)}
\begin{document}
\maketitle

\noindent\emph{Source and licence.} Fenchel-Willmore-Chen Inequality under Lower Bounds on weighted intermediate Ricci curvature. Version arXiv:2608.01171v1 (CC BY 4.0), distributed by arXiv under the Creative Commons Attribution 4.0 International licence, http://creativecommons.org/licenses/by/4.0/, as recorded in the arXiv licence metadata for this exact version and shown on the abstract page https://arxiv.org/abs/2608.01171v1. Full licence notice: \texttt{licenses/arxiv-2608.01171-CC-BY-4.0.txt}.\par\smallskip

\noindent\textit{Editorial scope.} Counted unit: the article's prop lem:weighted-bishop-jacobi (source0.tex lines 506--535). The article's complete original proof of it is delivered with it (source0.tex lines 538--678, one \texttt{proof} environment of 141 lines). No mathematics is written, repaired or re-derived: the statement and the proof are the article's own lines, in the article's own notation, and no retained line is substituted or re-flowed.  Retained uncounted as the source-local context the proof needs: lines 275--290; lines 353--355.  Nothing was omitted from any retained span, so the package carries no omission record.  No retained line contains a citation, so the wrapper carries no bibliography and no bibliography entry.  No retained line contains a figure environment, an image-inclusion command or drawing code, so no image is delivered and the package declares no asset dependency.  Editorial formatting only, in addition: an AMS \texttt{amsart} preamble that restates the article's own declarations and macros used by the retained lines, namely the environments \texttt{prop} on separate counters, together with the standard packages those lines need.  The retained environments print in this document as Proposition 1 in order of appearance, whereas the article prints the same items with the numbering of its own full document.  The complete native source of the article as distributed in the arXiv e-print for this exact version is bundled unchanged as \texttt{sources/206603/source0.tex}.
        \begin{prop}[Weighted partial Hessian comparison]\label{prop-weighted hessian}
        Let $(M^n,g,e^{-f}d\operatorname{vol})$ be a complete smooth metric measure space.
        Let
        $r_q(y):=d(q,y)$ be the distance function from a fixed point $q \in M$.
        Assume that
        \begin{equation}\label{Curvature-assumption}
        \mathrm{Ric}_{\ell,f}^1 \geq \ell K e^{-\frac{4f}{n-1}}
        \end{equation}
        for some $ K \in \mathbb{R}$ and some $1\leq \ell\leq n-1$.
        Let $p \in M \setminus (\mathrm{cut}(q)\cup \{q\})$, and let $\sigma:[0,T] \to M$ be the unique unit-speed minimizing geodesic from $q$ to $p$, where $T = r_q(p)$. Let $$V_l \subset (\nabla r_q)_{p}^\perp \subset T_p M$$ be an $\ell$-dimensional subspace, and let $\{e_1, \ldots, e_\ell\}$ be an orthonormal basis of $V_l$. Denote by $E_i(t)$ the parallel vector field along $\sigma$ satisfying $E_i(T) = e_i$.
        Then for every $t \in (0,T]$, with the additional restriction
        $s(t)< \pi / \sqrt{K}$ when $K>0$, 
        we have
        $$\sum_{i=1}^\ell (\operatorname{Hess} r_q)_{\sigma(t)}(E_i(t), E_i(t)) \leq \ell e^{-\frac{2f(t)}{n-1}} \frac{\mathrm{cs}_K(s(t))}{\mathrm{sn}_K(s(t))} + \frac{\ell}{n-1} f'(t),$$
        where $f(t) : = f(\sigma(t))$, and $s(t)$ is defined in \eqref{def:s}.
        \end{prop}

\begin{equation}\label{def:s}
    s(t) = \int_0^{t} e^{-\frac{2f(\sigma(\tau))}{n-1}}d \tau.
\end{equation}

\begin{prop}[Comparison of the determinant of $\ell$-Jacobi fields for $\operatorname{Ric}^{1}_{\ell,f}$]
    \label{lem:weighted-bishop-jacobi}
Let $(M^n,g,e^{-f}d\operatorname{vol})$ be a complete smooth metric measure space.
Fix $q \in M$ and $v \in T_qM$ with $|v| =1$. Consider the unit-speed minimizing geodesic
$$
    \sigma(t)=\exp_q(tv),\qquad 0 \leq t < T_{\mathrm{cut}}(q,v).
$$
Let $1\leq \ell\leq n-1$, and let
$w_1,\ldots,w_\ell\in T_qM$ be linearly independent vectors perpendicular to $v$.
Define the Jacobi fields
$$
    Y_i(t):=(D\exp_q)_{tv}(t w_i),\qquad 1\leq i\leq \ell .
$$
I.e. $Y_i$ is the Jacobi field with $Y_i(0) =0, \ Y'_i(0) = w_i$. 
Assume that
$$
    \operatorname{Ric}^{1}_{\ell,f}
    \geq
    \ell K e^{-\frac{4f}{n-1}}.
$$
Then, for every $0<t< T_{\mathrm{cut}}(q,v)$,
$$
\frac{|Y_1(t)\wedge\cdots\wedge Y_\ell(t)|}
{|w_1\wedge\cdots\wedge w_\ell|}
\leq
e^{\frac{\ell}{n-1}\left(f(\sigma(0))+f(\sigma(t))\right)}
\operatorname{sn}_K^\ell(s(t)) ,
$$
where $s(t)$ is defined in $\eqref{def:s}$.
\end{prop}

\begin{proof}
Let $L_t:v^\perp\to \sigma'(t)^\perp$ be a map defined by
$$
L_t (w)=(D\exp_q)_{tv}(t w).
$$
Then $Y_i(t)=L_t(w_i)$. Since $q$ has no conjugate points along $\sigma$,
$L_t$ is nonsingular for $t \in (0, T_{\mathrm{cut}}(q,v))$. Define
$$
W(t):=
\operatorname{span}\{Y_1(t),\ldots,Y_\ell(t)\}
$$
and
$$
J_\ell(t):=|Y_1(t)\wedge\cdots\wedge Y_\ell(t)|.
$$
In particular, $J_\ell (t) >0 $ for $t \in (0, T_{\mathrm{cut}}(q,v))$.


Let
$$
S_t:=L_t' \circ L_t^{-1},
$$
where $L_t'w$ denotes the covariant derivative of $L_t(w)$ along $\sigma$.
Equivalently, $S_t$ is the shape operator of the geodesic
sphere centered at $q$, so that
$$
\langle S_t(X),X\rangle
=
\left(\operatorname{Hess} r_q\right)_{\sigma(t)}(X,X),
\qquad X\in \sigma'(t)^\perp.
$$

For each $t >0$, choose an orthonormal basis $\{E_1(t),\ldots,E_\ell(t)\}$  of $W(t)$.
We claim that 
$$
\frac{d}{dt}\log J_\ell(t)
=
\sum_{i=1}^\ell \langle S_t(E_i(t)),E_i(t)\rangle .
$$
Indeed, if $G(t)=(\langle Y_i(t),Y_j(t)\rangle)_{i,j=1}^\ell$, then
$J_\ell(t)^2=\det G(t)$. Hence
\begin{align*}
\frac{d}{dt}\log J_\ell(t)
&=
\frac12\operatorname{tr}\bigl(G(t)^{-1}G'(t)\bigr).
\end{align*}
Using $Y_i'(t)=S_tY_i(t)$, this trace is precisely
$\operatorname{tr}_{W(t)}S_t$.

Applying Proposition~\ref{prop-weighted hessian} to the $\ell$-plane
$W(t)\subset \sigma'(t)^\perp$ at each $t$, we obtain
$$
\frac{d}{dt}\log J_\ell(t)
\le
\ell e^{-\frac{2f(t)}{n-1}}
\frac{\operatorname{cs}_K(s(t))}{\operatorname{sn}_K(s(t))}
+
\frac{\ell}{n-1}f'(t),
$$
where $f(t) : = f(\sigma(t))$ and $s(t)$ is defined in equation~\eqref{def:s}.
Since
$$
s'(t)=e^{-\frac{2f(t)}{n-1}},
$$
we have
$$
\frac{d}{dt}\log \operatorname{sn}_K(s(t))
=
e^{-\frac{2f(t)}{n-1}}
\frac{\operatorname{cs}_K(s(t))}{\operatorname{sn}_K(s(t))}.
$$
Therefore
$$
\frac{d}{dt}
\left(
\log J_\ell(t)
-
\ell\log \operatorname{sn}_K(s(t))
-
\frac{\ell}{n-1}f(t)
\right)
\le 0.
$$
Hence, for $0<\varepsilon<t$,
$$
\frac{J_\ell(t)}
{e^{\frac{\ell}{n-1}f(t)}\operatorname{sn}_K^\ell(s(t))}
\le
\frac{J_\ell(\varepsilon)}
{e^{\frac{\ell}{n-1}f(\varepsilon)}
\operatorname{sn}_K^\ell(s(\varepsilon))}.
$$

Since
$$
J_\ell(\varepsilon)
=
\varepsilon^\ell |w_1\wedge\cdots\wedge w_\ell|
+
O(\varepsilon^{\ell+2}) \quad \text{as }\varepsilon \to 0^+
$$
and
$$
s(\varepsilon)
=
e^{-\frac{2f(q)}{n-1}}\varepsilon
+
O(\varepsilon^2) \quad \text{as } \varepsilon \to 0^+,
$$
we have the limit
\begin{align*}
\lim_{\varepsilon\to0^+}
\frac{J_\ell(\varepsilon)}
{e^{\frac{\ell}{n-1}f(\varepsilon)}
\operatorname{sn}_K(s(\varepsilon))^\ell}
&=
e^{\frac{\ell}{n-1}f(q)}
|w_1\wedge\cdots\wedge w_\ell|.
\end{align*}
Consequently,
$$
J_\ell(t)
\le
e^{\frac{\ell}{n-1}\left(f(q)+f(\sigma(t))\right)}
\operatorname{sn}_K^\ell(s(t))
|w_1\wedge\cdots\wedge w_\ell|.
$$
Since
$$
J_\ell(t)=|Y_1(t)\wedge\cdots\wedge Y_\ell(t)|,
$$
we obtain
$$
\frac{|Y_1(t)\wedge\cdots\wedge Y_\ell(t)|}
{|w_1\wedge\cdots\wedge w_\ell|}
\leq
e^{\frac{\ell}{n-1}\left(f(q)+f(\sigma(t))\right)}
\operatorname{sn}_K(s(t))^\ell .
$$
This proves the desired estimate.
\end{proof}

\end{document}
